\documentclass[usenames,dvipsnames,10pt,aspectratio=169]{beamer} 

\usepackage[utf8]{inputenc}
\usepackage{verbatim}
\usepackage{minted}
\usepackage{graphicx}
\usepackage{wrapfig}
\usepackage{geometry}
\usepackage{listings}
% \usepackage{showframe}
\usepackage{enumitem}
\usepackage{color, xcolor}
\usepackage[document]{ragged2e}
\usetheme{umu}

\usemintedstyle{monokai}

\usepackage{hyperref}
\hypersetup{
    colorlinks=true,
    linkcolor=ucugreyish,
    filecolor=ucured,
    urlcolor=ucublue,
}
\urlstyle{same}

%%% Some useful commands
% pdf-friendly newline in links
\newcommand{\pdfnewline}{\texorpdfstring{\newline}{ }} 
% Fill the vertical space in a slide (to put text at the bottom)
\newcommand{\framefill}{\vskip 0pt plus 1 filll}

%%% Enter additional packages below (or above, I can't stop you)! / Jesper
\renewcommand{\proofname}{\sffamily{Proof}}

% presentation template slides usage
% \framecard[color (not working)]{textbuf}
% \framesplit{Header}{picture}{textbuf}
% \framepic{image}{text}
% \lstinputlisting[language=Bash, style=codestyle]{code/namespace_ex.sh}

%%%%%%%%%%%%%%%%%%%%%%%%%%%%%%%%%%%%%%%%%%%%%%%%%%%%%%%%%%%%%%%%%%%%%%%%%%%%%%%%%%%%%
\title{Linux course}
\subtitle{Tools overview}
\date[\today]{\small\today}
\author[Morhunenko Mykola]{Morhunenko Mykola}
\institute{APPS@UCU}

\setlist[itemize, 1]{label=$\color{ucublue} \bullet$, leftmargin=-2mm}

\begin{document}

\begin{frame}
\titlepage
\end{frame}

\begin{frame}{\contentsname}
    \setbeamercolor{background canvas}{bg=ucugrey}
    \tableofcontents
\end{frame}

\begin{frame}
    \frametitle{Intro}
    \begin{itemize}
        \item In this presentation, we will overview some tools that are available on all Linux distributions
        \item All of them have a high barrier to entry as Linux itself, but when you are there - you won't imagine your life without that tools
        \item Example: it's not a one-day task to learn how to move around your system, but after few months of practice working with CLI, GUI for you will be as slow as a turtle is slow in comparison with a rabbit
    \end{itemize}
\end{frame}

\section{Vim}

\framepic{graphics/vim_w.png}{}
\framesplitc{Vi}{graphics/vim_meme.png}{
    \begin{itemize}
        \item So let's start with the text editor
        \item You have heard about \ex{vim}, haven't you?
        \item But let's start with \ex{vi}
        \item Vi is a part of POSIX
        \item It's totally CLI editor (forget that you have a mouse)
        \item There are shortcuts for everything
        \item If there is no, you can create them for yourself
        \item Every good enough 21'st century editor has an extension for a \ex{vi mode}
        \item But almost nobody uses it - only on some low memory and low power machines. So we move to vim
    \end{itemize}
}

\begin{frame}{Vim}
    \begin{itemize}
        \item \ex{Vim} stands for 'Vi IMproved'
        \item According to Linux Journal survey, 38\% (in average for 2009-2018) of respondents vote for vim as the best editor
        \item It has much more features, than \ex{vi}, including more commands, scriptable syntax highlighting and extensions, graphical interface (and a mouse support, but don't use it)
        \item As \ex{vi}, it has six modes - normal, visual, insert, command-line, select, and ex (yes, not only NORMAL and INSERT)
        \item Because of a huge community (38\% of world's best geeks) \ex{vim} became a powerful IDE with thousands of extensions (syntax highlight, autocompletion, spell checking, project tree etc)
        \item The most powerful tool of \ex{vim} is inside - its shortcuts. You can make your work dozens of times faster without a touchpad and a mous
    \end{itemize}
\end{frame}

\begin{frame}{NeoVim}
    \begin{itemize}
        \item Neovim is just a fork of Vim with some Python extensions 
        \item And cool logo =)
        \item Also Neovim is a community-driven text editor, while Vim is a project of only one person - \ex{Bram Moolenaar}
        \item One 'expert' on reddit wrote that:
        
                \hspace*{1cm} "Neovim exists to convince Bram to push new features to Vim"
            
            And I mostly agree with him.
    \end{itemize}    
\end{frame}


\section{Tmux}
{ % all template changes are local to this group.
    \setbeamertemplate{navigation symbols}{}
    \begin{frame}<article:0>[plain]
        \begin{tikzpicture}[remember picture,overlay]
            \node[at=(current page.center)] {
                \includegraphics[keepaspectratio,width=\paperwidth/2]{graphics/Tmux_logo.png}
            };
        \end{tikzpicture}
     \end{frame}
}

\begin{frame}{Tmux}
    \begin{itemize}
        \item \ex{TMUX} stands for Terminal MUltipleXer
        \item There are some other (screen, Konsole, etc.), but they are not so good as TMUX
        \item Why do we need it?
        \item As you continue your practice in CLI, you can notice that it's not enough to have only one terminal window
        \item With this much multitasking going on, we want to have more terminals. So people create a \ex{terminal multiplexor}
        \item What TMUX can do?
        \begin{itemize}
            \item Not only split and stack tab but also make tabs
            \item Continue running programs in the background
            \item With extensions you can write layout files in \ex{.yml} format
            \item Search through terminal output and move around with Vim shortcuts
            \item Other interesting stuff
        \end{itemize}
    \end{itemize}
\end{frame}


\framesplitc{Tmux}{graphics/tmux_explain.png}{
    \begin{itemize}
        \item A Tmux Session with two tmux tabs with multiple tmux panes within each
        \item As \ex{vim}, tmux has modes - view and command (Ctrl+b be default)
        \item Every pane has three modes - view, choose and copy
        \item To enter a \ex{copy mode} - \ex{Ctrl+b [}
        \item It allows you to use vim keys for moving around and copying text
        \item For more information see \ex{man tmux} or \href{https://man7.org/linux/man-pages/man1/tmux.1.html}{linux man page}
    \end{itemize}
}

\begin{frame}{For the very beginning:}
        \hspace{-1cm}
        \includegraphics[width=\textwidth + 1.8cm, height=\textheight - 1cm]{graphics/tmux.png}
\end{frame}


\section{Ranger}
{ % all template changes are local to this group.
    \setbeamertemplate{navigation symbols}{}
    \begin{frame}<article:0>[plain]
        \begin{tikzpicture}[remember picture,overlay]
            \node[at=(current page.center)] {
                \includegraphics[keepaspectratio,width=\paperwidth/4]{graphics/ra.png}
            };
        \end{tikzpicture}
     \end{frame}
}
\begin{frame}{Ranger}
    
\end{frame}

\section{i3}
\framepic{graphics/i3wm.png}{i3wm}
\begin{frame}{i3}
    \begin{itemize}
        \item 
    \end{itemize}
\end{frame}

\section{Package managers}
\framepic{graphics/pacman.jpg}{Package 

managers}
\begin{frame}{Package managers}
    
\end{frame}



\section{Sources}
\framecard{Sources}
\begin{frame}{Sources}
    \begin{itemize}
        \item \href{https://www.linuxjournal.com/node/1203742}{Linux journal}
        \item \href{https://opensource.com/article/21/5/linux-terminal-multiplexer}{Termianl Multiplexers}
        \item \href{https://protechnotes.com/comprehensive-tmux-tutorial-for-beginners-with-a-cheat-sheet/}{Tmux tutorial}
        \item \href{https://man7.org/linux/man-pages/man1/tmux.1.html}{Tmux Linux man page}
    \end{itemize}
\end{frame}

\end{document}